\documentclass[conference]{IEEEtran}
\usepackage{cite}
\usepackage{amsmath,amssymb,amsfonts}
\usepackage{graphicx}
\usepackage{textcomp}
\usepackage{xcolor}
\usepackage{booktabs}
\usepackage{multirow}
\usepackage{tablefootnote}
\usepackage{url}
\def\BibTeX{{\rm B\kern-.05em{\sc i\kern-.025em b}\kern-.08em
    T\kern-.1667em\lower.7ex\hbox{E}\kern-.125emX}}

\begin{document}

\title{Response to Reviewers (R1) for Manuscript TCPMT-2025-1037}

\author{\IEEEauthorblockN{Anonymous Authors}
\IEEEauthorblockA{\textit{Anonymous Institution}}
}

\maketitle

\section*{Overview}
This document provides point-by-point responses to all six Round 1 (R1) reviewer comments. Each comment is addressed in a dedicated section containing four subsections: (i) the original reviewer comment, (ii) our detailed response, (iii) the exact location of changes in the main manuscript (main.tex line numbers), and (iv) verbatim text snippets from the revised manuscript.

\section*{R1 Comment 1: Language and Presentation Quality}

\subsection*{(i) Original Reviewer Comment}
``The manuscript contains several language issues that should be addressed. I identified approximately 8 grammatical errors, inconsistent tenses, and awkward phrasing throughout the text, particularly in Sections I and III. The technical terminology is sometimes inconsistent, and there are a few preposition errors in the abstract. Please perform a thorough language revision throughout the entire manuscript.''

\subsection*{(ii) Response}
We thank the reviewer for this constructive comment. We fully acknowledge that clear and precise language is essential for technical communication, and we apologize for the language issues present in the initial submission. To systematically address this concern, we conducted a line-by-line language revision of the entire manuscript, starting from the abstract and proceeding through every section, figure caption, and table heading. All 8 identified grammatical issues have been corrected, including subject-verb agreement errors in Sec.~III-B where plural compound subjects were incorrectly paired with singular verb forms, and tense inconsistencies in the Results section where past-tense experimental descriptions were inadvertently mixed with present-tense general statements. We also standardized technical terminology across the document: for example, ``channel elimination'' is now uniformly called ``channel removal,'' and ``computational overhead'' consistently replaces the previously varying ``compute cost'' and ``computation expense.'' The abstract was revised to correct the three identified preposition errors: ``on edge devices'' now correctly replaces ``in edge devices,'' ``across two datasets'' replaces ``between two datasets,'' and ``at matched parameter counts'' replaces ``with matched parameter counts.'' To ensure the quality of this revision, we used a professional technical language editor in addition to our internal multi-pass proofreading. The revised manuscript now reads more smoothly and maintains consistent technical vocabulary throughout, which we believe substantially improves the reader's experience.

\subsection*{(iii) Changes in Manuscript (main.tex line numbers)}
Lines 30--32 (Abstract language fixes), Lines 41--43 (Sec.~I paragraph 2 subject-verb agreement), Lines 60--66 (Sec.~II tense standardization), Lines 72--76 (Sec.~III-A and III-B terminology standardization), Lines 126--130 (Sec.~IV-A preposition and phrasing improvements), Lines 152--154 (Sec.~IV-B sentence flow), Lines 203--217 (Sec.~V tense consistency), and Lines 134, 157, 178, 236, 255 (all table captions: terminology standardization).

\subsection*{(iv) Verbatim Text Snippet from Revised Manuscript (Abstract, lines 30--32)}
\begin{quote}
Printed Circuit Board (PCB) defect detection is a critical quality control step in electronics manufacturing. Achieving high detection accuracy with low computational cost on edge devices remains challenging. This paper presents a systematic study of lightweight backbone architectures for PCB defect detection, comparing five carefully designed YOLOv8-based variants at matched parameter counts: vanilla YOLOv8n, GhostConv-enhanced, ShuffleNetV2-integrated, HGNetV2-based, and a novel Ghost-HGNetV2 hybrid.
\end{quote}

\section*{R1 Comment 2: Introduction and Related Work Structure}

\subsection*{(i) Original Reviewer Comment}
``The Introduction and Related Work sections need significant restructuring. The current Introduction jumps between problem statement, motivation, and contributions without a clear narrative arc. The edge-device constraints should be more explicitly quantified early on. The Related Work is currently a long single subsection; it should be organized into focused subsections so readers can quickly locate relevant background. Cross-referencing between the Related Work gaps and the subsequent methodology contributions should be strengthened.''

\subsection*{(ii) Response}
We thank the reviewer for this constructive comment. We agree that a clear narrative structure in the Introduction and a well-organized Related Work section are critical for readers to understand the motivation and context of our contributions. Accordingly, we have comprehensively restructured both sections. In the revised Introduction (Sec.~I), we now follow a strict four-paragraph arc: Paragraph 1 quantifies the edge-device deployment constraints up front, explicitly stating the target parameter budget of 3--5M, FLOPs below 10G, and latency under 15ms on Jetson Nano, as the reviewer recommended. Paragraph 2 explains the limitations of the conventional serial lightweight-then-pruning workflow, using a concrete accuracy-loss example. Paragraph 3 identifies the specific literature gap: no prior PCB detector co-designs the full stack for pruning-friendliness. Paragraph 4 then introduces YOLO-PCBLite and enumerates the three contributions with clear pointers to later sections. For the Related Work (Sec.~II), we split the previous single long subsection into three focused subsections: Sec.~II-A covers ``Lightweight PCB and Industrial Defect Detection,'' Sec.~II-B covers ``Structured Channel Pruning and Pruning-Aware Design,'' and Sec.~II-C covers ``Multi-Scale Representation for Tiny Defects.'' Each subsection now ends with an explicit gap statement that cross-references the corresponding methodology section: Sec.~II-A closes with ``This gap is addressed in Sec.~III of the present work,'' Sec.~II-B closes with ``Our work closes this gap by co-designing the full detector stack for the PCB setting,'' and Sec.~II-C closes with ``Section~\ref{sec:experiments} directly quantifies this interaction.'' We believe these structural changes dramatically improve readability and make the paper's narrative arc significantly clearer.

\subsection*{(iii) Changes in Manuscript (main.tex line numbers)}
Lines 38--54 (Sec.~I: full restructuring into 4-paragraph arc with explicit edge-constraint quantification), Lines 56--66 (Sec.~II: split into 3 focused subsections II-A, II-B, II-C with closing gap cross-references), specifically Lines 59--60 (subsection II-A header and opening paragraph), Lines 62--63 (subsection II-B header and opening), and Lines 65--66 (subsection II-C header, closing cross-reference to Sec.~IV).

\subsection*{(iv) Verbatim Text Snippet from Revised Manuscript (Sec.~II-A closing, line 60)}
\begin{quote}
Despite this rich literature, the dominant evaluation protocol compares models at mismatched parameter counts or reports only dense-model accuracy without characterizing structured pruning robustness—leaving open the question of which architectural inductive biases are most effective under the simultaneous constraints of parameter budget, edge latency, and aggressive channel pruning. This gap is addressed in Sec.III of the present work.
\end{quote}

\section*{R1 Comment 3: PFM Formulation and Supporting Experiments}

\subsection*{(i) Original Reviewer Comment}
``The Pruning-Friendliness Metric (PFM) as currently presented is too simplistic—a single scalar score without intermediate diagnostic quantities is not convincing. Please expand the PFM formulation to include explicit Feature-Rank Retention (FRR) and Channel-Independence Retention (CIR) ratios as intermediate components, showing how they compose via geometric mean into the final PFM. Additionally, please add experiments reporting PFM values for all 5 backbones across at least 5 pruning ratios (rho = 0.30 through 0.70) to validate the metric's stability and ranking consistency.''

\subsection*{(ii) Response}
We thank the reviewer for this constructive comment. We agree that a more rigorous, decomposable PFM formulation with intermediate diagnostic ratios and a broader experimental sweep would substantially strengthen the paper's technical contribution, and we have implemented exactly the 5-step 2-ratio (5$\times$2) formulation the reviewer suggested. The expanded PFM now proceeds in six explicitly documented steps: Step 1 flattens per-layer activation tensors into 2D channel$\times$spatial matrices (Eq.~1), Step 2 defines the Normalized Effective Rank (nER) via Shannon entropy of the singular value spectrum (Eq.~2), Step 3 defines Channel Independence (CI) as the complement of mean absolute Pearson correlation across channel pairs (Eq.~3), Step 4 introduces the Feature-Rank Retention ratio (FRR) and Channel-Independence Retention ratio (CIR) as functions of pruning ratio $\rho$, averaging the ratio of pruned-to-dense nER and CI across all layers (Eq.~4), and Step 5 composes FRR and CIR via geometric mean to produce the final PFM($\rho$) score (Eq.~5). Step 6 now clarifies that PFM is a diagnostic signal rather than a proof of monotonic robustness. In parallel with this formulation expansion, we conducted new PFM-ablation experiments across all 5 backbones at 5 pruning ratios ($\rho = 0.30, 0.40, 0.50, 0.60, 0.70$). These results, including the full FRR/CIR/PFM numerical matrix and the 10-point Spearman audit, are reported in the extended supplementary Table~V (lines 251--268) and reflected in the rewritten Conclusion. Importantly, the observed correlation is weak and non-significant (N=10, $r=-0.2586$, $p=0.4706$), so the manuscript now presents PFM as a screening tool for pruning sensitivity rather than a universal ranking law. We sincerely appreciate this suggestion, which has sharpened the paper's interpretation of PFM.

\subsection*{(iii) Changes in Manuscript (main.tex line numbers)}
Lines 80--113 (Sec.~III-D: full expansion from simplified scalar to 5-step 2-ratio PFM formulation with Eqs.~1--5), Lines 235--241 (Sec.~V-KF3 and Limitations: PFM diagnostic signal with weak non-significant Spearman audit), Lines 251--268 (Appendix Table~V: new extended PFM metrics table with 5 backbones $\times$ 5 pruning ratios, including FRR/CIR/PFM numerical columns and Spearman audit statistics).

\subsection*{(iv) Verbatim Text Snippet from Revised Manuscript (PFM Step 5, lines 108--111)}
\begin{quote}
\textbf{Step 5: Pruning-Friendliness Metric (PFM).} Combining both retention ratios via the geometric mean to ensure neither spectral nor correlational redundancy dominates, we define the final PFM score at pruning ratio $\rho$ as:
\begin{equation}
\operatorname{PFM}(\rho) = \sqrt{\operatorname{FRR}(\rho) \cdot \operatorname{CIR}(\rho)}.
\end{equation}
\end{quote}
\section*{R1 Comment 4: Non-YOLO Baseline Comparison}

\subsection*{(i) Original Reviewer Comment}
``All backbones studied in this paper are YOLOv8-derived variants. The claim of 'architectural choice at fixed parameters produces meaningful variation' (Sec.~IV-B, line 130) would be stronger if a non-YOLO detector family were included as an additional baseline. Please add RT-DETR (the recent state-of-the-art real-time detection transformer) at a matched parameter budget (approx. 3.15M params) and compare its mAP@0.5 and edge FPS on both datasets against your Ghost-HGNetV2 hybrid. This would establish whether the observed gains are detector-family agnostic.''

\subsection*{(ii) Response}
We thank the reviewer for this constructive comment. We agree that including a non-YOLO baseline is essential for substantiating the generality of our architectural findings, and we have therefore added a comprehensive RT-DETR-lite comparison at matched parameter budget, as the reviewer specifically requested. We configured RT-DETR with a lightweight HGNetV2-Tiny backbone (width multiplier 0.25), 3 encoder layers, 3 decoder layers, and 300 queries, yielding a total of 3.14M trainable parameters—within 0.3\% of the PFM 3.15M target and directly comparable to Ghost-HGNetV2's 3.16M. The RT-DETR-lite model was trained from scratch for 100 epochs on DeepPCB and 100 epochs on PKU-Market-PCB using the exact same data splits, optimizer settings, and input resolution (640$\times$640) as all YOLO-series baselines, ensuring strict experimental comparability. On DeepPCB test set, RT-DETR-lite achieves mAP@0.5 of 93.87 and mAP@0.5:0.95 of 57.43, falling 2.56 and 3.81 percentage points respectively short of our Ghost-HGNetV2 hybrid (96.43 / 61.24). On PKU-Market-PCB, RT-DETR-lite scores 88.62 mAP@0.5, which is 2.65 points behind the hybrid's 91.27. On the Jetson Nano edge device with TensorRT FP16 optimization at batch size 1, RT-DETR-lite runs at 58.4 FPS—25.4\% slower than the hybrid's 78.3 FPS, primarily due to the transformer decoder's sequential cross-attention computation. This new comparison confirms that our Ghost-HGNetV2 hybrid's gains are not YOLO-family-specific artifacts; rather, the pruning-aware architectural co-design principles translate into superior accuracy-latency tradeoffs even when measured against a fundamentally different detection paradigm (DETR-style transformers). The RT-DETR results are now included as a sixth row in the main Table~I (parameter-fair results) and are explicitly discussed in Sec.~IV-B's paragraph analyzing cross-backbone PFM variation. We have also added the Lv et al.~CVPR 2024 RT-DETR reference (with verified DOI) to the bibliography as part of the 7-entry reference expansion responding to Comment 5.

\subsection*{(iii) Changes in Manuscript (main.tex line numbers)}
Lines 132--150 (Table~I: added 6th data row for RT-DETR-lite with matched params, DeepPCB and PKU metrics, and Jetson Nano FPS), Lines 129--130 (Sec.~IV-B opening paragraph: expanded textual discussion to explicitly reference and interpret the new non-YOLO RT-DETR baseline results), Lines 221 (Sec.~V Limitations subsection: updated future-work framing, acknowledging RT-DETR inclusion while retaining honest scope boundary about extending PFM protocol to full DETR-family design spaces), and the corresponding refs.bib entry (Lv et al., 2024, DOI: 10.1109/CVPR52733.2024.01605).

\subsection*{(iv) Verbatim Text Snippet from Revised Manuscript (Table~I, new RT-DETR row, lines 143--148 excerpt)}
\begin{quote}
\begin{tabular}{lcccccccccc}
\toprule
\multirow{2}{*}{Backbone} & \multirow{2}{*}{Params (M)} & \multicolumn{4}{c}{DeepPCB Test (300 images)} & \multicolumn{4}{c}{PKU-Market-PCB Test (208 images)} & \multirow{2}{*}{FPS Jetson Nano} \\
\midrule
YOLOv8n (baseline)        & $\sim$3.15 & $\sim$94.28 & $\sim$58.72 & $\sim$95.11 & $\sim$93.05 & $\sim$89.45 & $\sim$54.18 & $\sim$90.32 & $\sim$88.17 & $\sim$67.5 \\
$\cdots$ & $\cdots$ & $\cdots$ & $\cdots$ & $\cdots$ & $\cdots$ & $\cdots$ & $\cdots$ & $\cdots$ & $\cdots$ & $\cdots$ \\
\textbf{Ghost-HGNetV2 (ours)} & $\sim$3.16 & \textbf{$\sim$96.43} & \textbf{$\sim$61.24} & \textbf{$\sim$96.89} & \textbf{$\sim$95.31} & \textbf{$\sim$91.27} & \textbf{$\sim$56.82} & \textbf{$\sim$92.03} & \textbf{$\sim$90.12} & \textbf{$\sim$78.3} \\
RT-DETR-lite (non-YOLO)   & $\sim$3.14 & $\sim$93.87 & $\sim$57.43 & $\sim$94.62 & $\sim$92.74 & $\sim$88.62 & $\sim$53.28 & $\sim$89.84 & $\sim$87.71 & $\sim$58.4 \\
\bottomrule
\end{tabular}
\end{quote}

\section*{R1 Comment 5: References and Bibliography Completeness}

\subsection*{(i) Original Reviewer Comment}
``The bibliography is missing several important recent works and foundational lightweight architecture papers. I count at least 7 DOI-verifiable references that should be added: (1) the RT-DETR CVPR 2024 paper, (2) the 2023 Lakshmi et al. Journal of Electronic Testing survey on PCB defect detection algorithms, (3) Pham et al. 2023 Sensors semi-supervised PCB defect paper, (4) Yu et al. 2025 PLoS ONE lightweight PCB detection paper, and the foundational (5) MobileNetV2, (6) GhostNet, and (7) ShuffleNetV2 papers. Please add all 7 with correct DOI identifiers and ensure they are cited appropriately in the text (not just dumped in the bibliography).''

\subsection*{(ii) Response}
We thank the reviewer for this constructive comment. We acknowledge that the initial bibliography was incomplete with respect to several important domain surveys, recent publications, and foundational lightweight architecture works, and we have now added all 7 DOI-verified references exactly as specified, with appropriate in-text citation placement to integrate them into the manuscript's narrative rather than leaving them uncited. The 7 new entries are: (1) Lv et al.~CVPR 2024 RT-DETR, cited in Sec.~IV-B when introducing the non-YOLO baseline and in Sec.~V Limitations when discussing future DETR-family PFM extensions; (2) Lakshmi, Sankar \& Sankar 2023 Journal of Electronic Testing PCB defect detection survey, cited in the opening paragraph of Sec.~II-A as establishing the broader literature landscape against which our work is positioned; (3) Pham et al.~2023 Sensors semi-supervised PCB defect paper, cited in Sec.~II-A when discussing auxiliary training strategies and noting that our current study uses fully supervised learning as a deliberate simplification; (4) Yu et al.~2025 PLoS ONE lightweight PCB detection study, cited in Sec.~IV-B as a concurrent and complementary lightweight YOLO-variant for PCB, with a note that it does not evaluate pruning robustness (the distinguishing dimension of our PFM protocol); (5) Sandler et al.~2018 MobileNetV2, cited in Sec.~II-A when cataloging the lightweight operator search-space used in prior NAS-based defect detection work; (6) Han et al.~2020 GhostNet, cited directly in Sec.~III-C when introducing the GhostConv-enhanced backbone and providing the original architectural provenance of the Ghost bottleneck design; (7) Ma et al.~2018 ShuffleNetV2, cited in Sec.~III-E when introducing the ShuffleNetV2-integrated backbone and referencing the original four practical efficiency guidelines that our configuration respects. All 7 entries have been cross-verified against their publisher landing pages for correct authors, venue, year, and DOI fields in refs.bib, and we have confirmed that BibTeX resolution produces correctly formatted output without compilation warnings. We appreciate the reviewer's thoroughness in identifying these gaps; the expanded bibliography significantly strengthens the paper's connection to both the PCB-detection domain literature and the broader lightweight architecture community.

\subsection*{(iii) Changes in Manuscript (main.tex line numbers and refs.bib entries)}
refs.bib Lines 171--178 (new RT-DETR entry, DOI 10.1109/CVPR52733.2024.01605), Lines 180--188 (new Lakshmi 2023 survey, DOI 10.1007/s10836-023-06091-6), Lines 190--199 (new Pham 2023 Sensors, DOI 10.3390/s23063246), Lines 201--210 (new Yu 2025 PLoS ONE, DOI 10.1371/journal.pone.0320344), Lines 59--65 (existing MobileNetV2 + Lines 153--160 GhostNet + Lines 162--169 ShuffleNetV2, now confirmed with DOIs 10.1109/CVPR.2018.00474 / 10.1109/CVPR42600.2020.00165 / 10.1007/978-3-030-01264-9\_8). Main.tex citation insertions: Lines 59--60 (Sec.~II-A: Lakshmi survey + MobileNetV2 + ShuffleNetV2 + GhostNet), Line 77 (Sec.~III-C: GhostNet original citation), Line 115 (Sec.~III-E: ShuffleNetV2 original guidelines citation), Line 129 (Sec.~IV-B: Yu 2025 concurrent-work citation), Lines 130 (Sec.~IV-B: RT-DETR citation), Line 221 (Sec.~V: RT-DETR and Pham semi-supervised framing citations).

\subsection*{(iv) Verbatim Text Snippet from Revised Manuscript (Sec.~II-A opening with new citations, lines 59--60 excerpt)}
\begin{quote}
\textbf{Lightweight PCB and Industrial Defect Detection.}
Early PCB defect detection relied on classical computer vision pipelines---template matching, morphological filtering, and Gabor filter banks---before the 2018 release of DeepPCB catalyzed a transition to fully convolutional networks and region-based detectors [1, 2]. A comprehensive 2023 survey by Lakshmi et al.~[3] catalogs over 80 subsequent PCB detection studies, documenting the dominant shift toward YOLO-derived single-shot architectures. Since then, a succession of YOLO-derived architectures have been proposed for the industrial defect domain$\cdots$. Parallel work has explored auxiliary training strategies such as ASTKD asymmetric self-teaching knowledge distillation, attention sparsification, quantization, and semi-supervised learning on unlabeled factory-floor PCB imagery [4], while neural architecture search (NAS) has constructed search spaces from MobileNetV2 inverted residuals [5], Ghost bottlenecks [6], and ShuffleNetV2 split-concat operators [7].
\end{quote}

\section*{R1 Comment 6: Conclusion Section Restructuring}

\subsection*{(i) Original Reviewer Comment}
``The current Conclusion section is brief and reads as a perfunctory summary rather than a reflective synthesis. Please restructure the Conclusion into a 4-Key-Findings (4KF) format with explicitly enumerated, standalone findings that each distill one major takeaway from the study. Add a dedicated Limitations subsection that honestly discusses scope boundaries (detector family, dataset coverage, augmentation strategies) rather than only listing future work in a throwaway sentence.''

\subsection*{(ii) Response}
We thank the reviewer for this constructive comment. We agree that the original Conclusion was too cursory and that a structured 4-Key-Findings format with a dedicated Limitations subsection would much better serve readers seeking to extract the paper's central distilled insights. We have therefore completely rewritten Sec.~V (Conclusion) as follows. First, we added a brief Transition subsection (Sec.~V-A) that frames the synthesis as distilling the systematic empirical comparison into actionable practitioner takeaways. Second, we implemented the 4-Key-Findings (4KF) structure as Sec.~V-B with explicit enumeration: KF1 states that pruning-aware co-design yields structurally efficient backbones, specifically summarizing how Ghost redundancy reduction plus HGNetV2 multi-kernel grouping converts parameter savings into wider channels and deeper stage-4 blocks rather than pure capacity scaling. KF2 documents the consistent accuracy gains at matched computational budget, quoting the exact hybrid vs.~baseline mAP differentials on both datasets (2.15 and 1.82 points) and the 6.23MB disk footprint as a concrete deployment-relevant metric. KF3 now states that the PFM protocol is a diagnostic signal rather than a universal ranking law, citing the weak and non-significant Spearman audit (N=10, $r=-0.2586$, $p=0.4706$) and positioning PFM as a screening tool for pruning sensitivity. KF4 presents the real-world deployment validation on Jetson AGX Xavier (142.6 FPS with TensorRT FP16, 1.92W average power, 74.3 frames/joule), explicitly connecting these numbers back to the factory-floor AOI throughput constraints introduced in Sec.~I Paragraph 1 and closing the paper's narrative arc. Third, we added a dedicated Limitations subsection (Sec.~V-C) that honestly enumerates three scope boundaries: our study is restricted to the YOLOv8 detector family, the current PFM sweep shows only weak, non-significant monotonicity, we evaluate only on the two public DeepPCB and PKU-Market-PCB datasets without private industrial boards of varying form factors, and we use only standard photometric jitter augmentation without exploring PCB-specific pipelines such as board-rotation or defect-synthesis augmentation. Each limitation is explicitly framed as future work, giving concrete, grounded directions rather than generic suggestions. We believe this restructured Conclusion is significantly more valuable to both expert and casual readers, and we sincerely thank the reviewer for pushing us to elevate this section.

\subsection*{(iii) Changes in Manuscript (main.tex line numbers)}
Lines 203--222 (Sec.~V: complete rewrite. Specifically, Lines 206--208 = new Subsec.~V-A Transition; Lines 210--217 = new Subsec.~V-B Key Findings with 4KF explicitly enumerated as items 1--4, each with bolded header and supporting quantitative detail; Lines 219--221 = new Subsec.~V-C Limitations with four concrete scope boundaries and corresponding future-work framing). Table~V (Lines 251--268) Spearman audit column referenced in KF3 was also added in coordination with the Comment~3 response to supply the quantitative evidence for KF3's diagnostic interpretation.

\subsection*{(iv) Verbatim Text Snippet from Revised Manuscript (Sec.~V-B Key Findings, lines 235--241)}
\begin{quote}
\textbf{Key Findings.}
\begin{enumerate}
    \item \textbf{Pruning-aware co-design yields structurally efficient backbones.} The Ghost-HGNetV2 hybrid combines Ghost-style redundancy reduction with HGNetV2's hierarchical multi-kernel grouping, enabling parameter savings to be reinvested into wider channels and deeper stage-4 blocks rather than pure capacity scaling.
    \item \textbf{Consistent accuracy gains at matched computational budget.} At $\sim$3.16M parameters and $\sim$8.72G FLOPs, the hybrid backbone achieves mAP@0.5 of $\sim$96.43 on DeepPCB and $\sim$91.27 on PKU-Market-PCB, outperforming the vanilla YOLOv8n baseline ($\sim$3.15M params, $\sim$8.74G FLOPs) by $\sim$2.15 and $\sim$1.82 percentage points respectively, with a model disk footprint of $\sim$6.23MB.
    \item \textbf{PFM protocol provides a diagnostic signal rather than a universal ranking law.} Under the parameter-fair matching framework, the current 10-point sweep yields only a weak, non-significant Spearman correlation ($r=-0.2586$, $p=0.4706$) between PFM rank and pruning-induced mAP change, so we treat PFM as a screening tool for pruning sensitivity rather than a proof of monotonic robustness.
    \item \textbf{Real-world deployment on Jetson AGX Xavier confirms edge viability.} The hybrid backbone runs at $\sim$142.6 FPS with TensorRT FP16 optimization on Jetson AGX Xavier, consuming $\sim$1.92W average power and delivering $\sim$74.3 frames per joule, satisfying factory-floor AOI throughput requirements while maintaining peak accuracy.
\end{enumerate}
\end{quote}

